\scanpage{332}\mspage{248}
% l. 1–4
$2^\circ)$ Les limites inductives de type $I_{\unsure{\alpha'|\alpha}}$ sont représentables dans les catégories $\unsure{\mathcal E|\mathscr E}_j$, et les foncteurs \unsure{cochgt|cochgt.} de base $\unsure{\alpha|\lambda}_*:\mathcal E_j\longrightarrow\mathcal E_{j'}$ y commutent.
% l. 4–11
[ C'est le cas p. ex. si les limites inductives filtrantes \uline{dénombrables} sont repr. dans les $\mathcal E_j$, et si les $\alpha_*$ y commutent ; ou s'il existe un cardinal $\unsure{\pi|\kappa}_0\leq\pi$, et un cardinal $\pi'_0>\pi_0$, tel que les limites inductives grandes devant $\pi_0$ et de cardinal $\leq\pi'_0$ soient repr. dans les $\mathcal E_j$, et que les $\alpha_*$ y commutent$\unsure{]|)|{}}$ ].

% l. 12–16
\margin{$\unsure{]\!]|]|\parallel}$}\unsure{Sous|Dans} ces conditions, le foncteur canonique fidèle
\[
 \underline{\Gamma}(\mathcal E/J)\longrightarrow\prod_{j\in\operatorname{Ob}J}\mathcal E_j \tag*{(4.10.3.1)}
\]
est cofinal, et la catégorie source est grde devant $\pi$.

% l. 17–19
Il est \inserted{$\pm$} immédiat que les hyp. de 4.10.1 impliquent que celles de 4.10.3 sont satisfaites pour $\mathcal E:\mathscr D_{\unsure{\varphi|\phi}}\longrightarrow J$ (4.10.5), Pour
% archive sheet number: 331
